\section{Reading map}\label{sec:overview}

Here is the argument in one pass.  A syndrome \(f\) determines a linear
system \(W_f\) of binary forms through its two-row Hankel matrix.  The
Hankel criterion says that a completely split squarefree form in \(W_f\)
is a witness that \(f\) lies in the span of the required number of distinct
normal-rational-curve points.  Thus the geometric problem is to find such a
form; only after the separate covering-radius gate does its absence mean that
\(f\) is a deep-hole syndrome.

At redundancy five, \(W_f\) is a pencil of cubics.  Its off-diagonal fibre
square records ordered pairs of distinct roots lying in one pencil member.
The exact count in Proposition~\ref{prop:r5-count} turns rational
points on that genus-one curve into split cubics and makes R5 the terminal
model for every later level.

At higher redundancy, choose a root \(\lambda\) as a marker and contract the
syndrome.  A split witness in the lower system lifts after multiplication by
\(\lambda\), and it stays squarefree precisely when it avoids that marker.
As \(\lambda\) varies, the contracted syndromes trace a polar line.  There
are then two branches.  If the line is transverse to the lower bad locus,
only finitely many markers and points on the lower splitting cover must be
discarded; a rational-point bound supplies a witness.  If the line is
contained, the component calculation places the original syndrome in the
persistent catalecticant locus or a characteristic-dependent Lucas carrier.
Figure~\ref{fig:polar-flow} displays this single recursion.

\begin{figure}[H]
\centering
\[
\begin{array}{rcl}
\text{coding}
 &:& [f]\longmapsto W_f\supset g_{\mathrm{split}}
      \longmapsto [f]\text{ is shallow},\\[4pt]
\text{recursion}
 &:& [f]\xrightarrow{\ \text{contract at }\lambda\ }
      \text{lower splitting problem},\\[2pt]
&&\begin{array}{rcl}
\text{contained}&\longmapsto&\text{persistent or Lucas carrier},\\
\text{transverse}&\longmapsto&
 \text{lower witness }g,\ \lambda\nmid g
 \longmapsto \lambda g\in W_f .
\end{array}
\end{array}
\]
\caption{The conceptual route through the paper.  The first line changes
from coding language to a splitting problem.  Marked contraction then
separates the contained branch, classified by the carrier theorem, from the
transverse branch, where finite avoidance and point counting produce the
lower witness \(g\).}
\label{fig:polar-flow}
\end{figure}

The proof has one recursive spine and six evaluated levels.
Table~\ref{tab:dependency-map} shows what changes from one level to the next;
its deletion budget is the number of points excluded from the terminal cover.
Table~\ref{tab:classification-status} separately records which geometric
classifications promote to code deep-hole classifications.

\begin{table}[H]
\centering
\small
\caption{Results and mechanisms.  The threshold column records the first
field range in which the stated geometric argument closes; bounded
certificates handle only the listed lower fields.}
\label{tab:dependency-map}
\setlength{\tabcolsep}{3pt}
\begin{tabular}{p{0.07\textwidth}p{0.24\textwidth}p{0.40\textwidth}
                p{0.22\textwidth}}
\toprule
Level & New obstruction & Main geometric mechanism & Closing gate\\
\midrule
R5 & cubic splitting & cubic pencil and off-diagonal \((2,2)\) fibre square;
exact split-witness count & branch correction at most \(12\), or \(6\) in
characteristic two; geometric threshold \(20\) and seven-field certificate\\
R6 & one retained marker & one contraction to the same R5 curve & deletion
budget \(19\); geometric threshold \(29\) and finite bridge\\
R7 & two retained markers & two coherent contractions to the same R5 curve &
deletion budget \(25\); geometric threshold \(37\) and finite bridge\\
R8 & three retained markers & the same integral \((2,2)\) terminal cover &
deletion budget \(30\); \(q\geq43\)\\
R9 & four markers and a characteristic-seven carrier & the terminal cover
off the carrier; a residual-quadratic genus-one slice on it & deletion budget
\(36\); \(q\geq53\), with the first relevant characteristic-seven field
handled geometrically\\
R10 & five markers and the first higher Lucas carrier & the terminal cover
in odd characteristic, an ordered-Hessian slice in characteristic two,
and a final-pair Artin--Schreier criterion on the carrier & deletion
budget \(42\); odd \(q\geq59\), even \(q\geq64\), plus the stated binary
carrier certificates\\
\bottomrule
\end{tabular}
\end{table}

\begin{table}[H]
\centering
\small
\caption{Classification status.  A split-free equality becomes a code
deep-hole equality only after the radius gate.}
\label{tab:classification-status}
\setlength{\tabcolsep}{3pt}
\begin{tabular}{p{0.07\textwidth}p{0.15\textwidth}
                p{0.29\textwidth}p{0.29\textwidth}
                p{0.12\textwidth}}
\toprule
Red. & Field range & Split-free geometry & Deep holes and radius
& Gap\\
\midrule
\(5\) & \(q\geq7\) & Complete & Complete; \(\rho=4\) & None\\
\(6\) & \(q\geq7\) & Complete & Complete; \(\rho=5\) & None\\
\(7\) & \(q\geq7\) & Complete & Complete for \(q\geq11\) & Radius at
\(7,8,9\)\\
\(8\) & \(q\geq43\) & Persistent-only & Complete; \(\rho=7\) & Below range\\
\(9\) & \(q\geq53\) & Persistent-only & Complete; \(\rho=8\) & Below range\\
\(10\) & \(q\geq59\) & Persistent-only & Complete; \(\rho=9\) & Below range\\
\(r\geq6\) & \(q\geq Q_r\) & Contained in persistent plus maximal Lucas;
persistent-only if \(p>r-1\), conditional on the intermediate lower
packages & Complete under that hypothesis if \(p>r-1\) & Stagewise packages;
Lucas arithmetic in small characteristic\\
\bottomrule
\end{tabular}
\end{table}

Here \(Q_r=6r-15+\lfloor2\sqrt{6r-17}\rfloor\).  The R7 finite record
is exposed only as a four-field exception delta beyond the uniform
\(T/T^6\) spine: the full fourteen-field ledger remains in the supplement.
